\chapter{Sialadenitis}

\section*{Definition}
Sialadenitis is the inflammation of salivary gland tissue, most commonly affecting the parotid or submandibular glands, characterized by painful swelling, tenderness, and purulent saliva from the duct orifice; may be acute (bacterial or viral) or chronic.

\section*{What Happens in Your Body}
Reduced salivary flow from dehydration, ductal blockage, or systemic illness allows retrograde migration of oral bacteria into the salivary parenchyma. Bacterial proliferation within the gland triggers neutrophilic infiltration, swelling, and collection of pus formation if untreated.

\section*{Causes}
\begin{itemize}
  \item \textbf{Bacterial (acute suppurative):} Staphylococcus aureus (most common), Streptococcus viridans, Haemophilus influenzae, anaerobes
  \item \textbf{Viral:} Mumps (paramyxovirus -- parotitis), Coxsackie, influenza, EBV, CMV, HIV
  \item \textbf{Ductal blockage:} Sialolithiasis (salivary duct stone) causing secondary infection
  \item \textbf{Autoimmune:} Sjogren syndrome (chronic, recurrent), IgG4-related disease
  \item \textbf{Radiation-induced:} Post-radiotherapy dry mouth and chronic sialadenitis
  \item \textbf{Medication-related:} Anticholinergics, diuretics, antihistamines (reduced salivary flow)
\end{itemize}

\section*{When to See a Doctor}
See your doctor if:
\begin{itemize}
  \item Acute swelling and pain over a salivary gland with fever and feeling unwell
  \item Purulent saliva expressed from the duct orifice on massage
  \item Trismus (difficulty opening the mouth) from parotid involvement
  \item Palpable stone along the duct (Wharton or Stensen duct)
  \item Recurrent or bilateral gland swelling suggesting autoimmune causes
  \item Rapid growth or firm, fixed mass (suspect neoplasm instead of infection)
\end{itemize}

\section*{Related Symptoms}
\begin{itemize}
  \item Gland swelling (parotid: preauricular; submandibular: submandibular triangle)
  \item Pain exacerbated by eating (sialodochitis)
  \item Trismus and difficulty swallowing
  \item Purulent or turbid saliva from the duct
  \item Fever, chills, and feeling unwell
\end{itemize}
\textbf{Important:} Bilateral parotid swelling in a febrile, unvaccinated child is mumps until proven otherwise; acute unilateral parotitis in a hospitalized, elderly, dehydrated patient is typically bacterial sialadenitis.

\section*{Self-Care / Home Management}
\begin{itemize}
  \item Oral hydration to stimulate salivary flow
  \item Warm compresses over the affected gland
  \item Sialogogues: tart candies, lemon drops, or sour citrus to promote saliva production
  \item Massage of the gland from posterior to anterior toward the duct orifice
  \item NSAIDs (ibuprofen) for pain and inflammation
  \item Good oral hygiene and regular dental care to reduce oral bacterial load
\end{itemize}

\section*{How Your Doctor Will Evaluate You}
\begin{itemize}
  \item Palpation of the gland and bimanual palpation of the duct for stones
  \item Expression of saliva from the duct orifice for inspection and culture
  \item Ultrasonography of the salivary gland to detect stones, collection of pus, or mass
  \item CT with contrast (sialolithiasis protocol) for deep lobe parotid or submandibular stones
  \item Sialography (MRI or conventional) for ductal anatomy and strictures
  \item Autoimmune serology (anti-Ro/SS-A, anti-La/SS-B, IgG4) for chronic or recurrent cases
\end{itemize}

\section*{Treatment Options}
\begin{itemize}
  \item \textbf{Acute bacterial:} IV or oral antibiotics (amoxicillin-clavulanate or clindamycin) for 7--10 days
  \item \textbf{Viral:} Supportive care (hydration, analgesia, fever-reducing medications); no antiviral indicated
  \item \textbf{collection of pus:} Incision and drainage with ultrasound or CT guidance
  \item \textbf{Sialolithiasis:} Sialendoscopy with stone retrieval, extracorporeal shock wave lithotripsy (ESWL), or gland excision
  \item \textbf{Chronic/recurrent:} Duct dilation, sialendoscopy with steroid irrigation
  \item \textbf{Autoimmune:} Oral hydroxychloroquine, pilocarpine for dry mouth, corticosteroid for acute flares
  \item \textbf{Radiation-induced:} Pilocarpine or cevimeline for dry mouth and sialadenitis prevention
\end{itemize}

\section*{References}
\begin{thebibliography}{99}
\bibitem{sialadenitis:parotid1} Wilson KF, Meier JD, Ward PD. ``Salivary gland disorders.'' Am Fam Physician. 2014;89(11):882-888.
\bibitem{sialadenitis:parotid2} Carlson ER, Ord RA. ``Salivary gland diseases.'' In: Peterson\'s Principles of Oral and Maxillofacial Surgery. 4th ed. Springer; 2022.
\bibitem{sialadenitis:parotid3} Schaitkin BM, Eisenman DS. ``Sialadenitis.'' In: Cummings Otolaryngology. 7th ed. Elsevier; 2021.
\bibitem{sialadenitis:parotid4} Albrecht M, Yanci KA. ``Sialadenitis.'' In: StatPearls. StatPearls Publishing; 2023.
\bibitem{sialadenitis:parotid5} Baum BJ, Fox PC. ``Salivary gland disease.'' In: Cecil Textbook of Medicine. 26th ed. Elsevier; 2020.
\bibitem{sialadenitis:parotid6} Walvekar RR, Andrade H, Simental AA. ``Salivary gland diseases in adults.'' In: Flint PW, et al. Cummings Otolaryngology. 7th ed. Elsevier; 2021.
\end{thebibliography}
